\documentclass[10pt,a4paper]{amsart}
\usepackage{amsmath,amssymb,amsthm,graphicx}
% Retained preamble for the Stacks Project chapter snapshot in this directory.
% The chapter's own preamble.tex is a copy of the shared Stacks Project preamble; the block
% below is sliced MECHANICALLY, line for line, from it. Nothing was added, removed or
% redefined; the wrapper class, the xr-hyper/hyperref wiring and the externaldocument list
% are omitted because the assembler supplies the document class.
%
% lines taken from preamble.tex: 3, 5, 9, 11, 12, 156, 157, 158, ... (51 lines total)
%
% package coverage restored from the chapter's own preamble (2026-10-13): xy, xy-2cell

\usepackage[all]{xy}
\xyoption{2cell}
\UseAllTwocells

\newenvironment{reference}{\comment}{\endcomment}
\newenvironment{slogan}{\comment}{\endcomment}
\newenvironment{history}{\comment}{\endcomment}
\theoremstyle{plain}
\newtheorem{theorem}[subsection]{Theorem}
\newtheorem{proposition}[subsection]{Proposition}
\newtheorem{lemma}[subsection]{Lemma}
\theoremstyle{definition}
\newtheorem{definition}[subsection]{Definition}
\newtheorem{example}[subsection]{Example}
\newtheorem{exercise}[subsection]{Exercise}
\newtheorem{situation}[subsection]{Situation}
\theoremstyle{remark}
\newtheorem{remark}[subsection]{Remark}
\newtheorem{remarks}[subsection]{Remarks}
\numberwithin{equation}{subsection}
\def\lim{\mathop{\mathrm{lim}}\nolimits}
\def\colim{\mathop{\mathrm{colim}}\nolimits}
\def\Spec{\mathop{\mathrm{Spec}}}
\def\Hom{\mathop{\mathrm{Hom}}\nolimits}
\def\Ext{\mathop{\mathrm{Ext}}\nolimits}
\def\SheafHom{\mathop{\mathcal{H}\!\mathit{om}}\nolimits}
\def\SheafExt{\mathop{\mathcal{E}\!\mathit{xt}}\nolimits}
\def\Sch{\mathit{Sch}}
\def\Mor{\mathop{\mathrm{Mor}}\nolimits}
\def\Ob{\mathop{\mathrm{Ob}}\nolimits}
\def\Sh{\mathop{\mathit{Sh}}\nolimits}
\def\NL{\mathop{N\!L}\nolimits}
\def\CH{\mathop{\mathrm{CH}}\nolimits}
\def\proetale{{pro\text{-}\acute{e}tale}}
\def\etale{{\acute{e}tale}}
\def\QCoh{\mathit{QCoh}}
\def\Ker{\mathop{\mathrm{Ker}}}
\def\Im{\mathop{\mathrm{Im}}}
\def\Coker{\mathop{\mathrm{Coker}}}
\def\Coim{\mathop{\mathrm{Coim}}}
\DeclareMathSymbol{\boxtimes}{\mathbin}{AMSa}{"02}
\def\QCohstack{\mathcal{QC}\!\mathit{oh}}
\def\Cohstack{\mathcal{C}\!\mathit{oh}}
\def\Spacesstack{\mathcal{S}\!\mathit{paces}}
\def\Quotfunctor{\mathrm{Quot}}
\def\Hilbfunctor{\mathrm{Hilb}}
\def\Curvesstack{\mathcal{C}\!\mathit{urves}}
\def\Polarizedstack{\mathcal{P}\!\mathit{olarized}}
\def\Complexesstack{\mathcal{C}\!\mathit{omplexes}}
\def\Pic{\mathop{\mathrm{Pic}}\nolimits}
\def\Picardstack{\mathcal{P}\!\mathit{ic}}
\def\Picardfunctor{\mathrm{Pic}}
\def\Deformationcategory{\mathcal{D}\!\mathit{ef}}
\usepackage{cleveref}
\crefname{originallemma}{original lemma}{original lemmas}
\Crefname{originallemma}{Original lemma}{Original lemmas}
\makeatletter
\newlabel{lemma-compare-cosk0}{{01GC}{source}{}{sourcepointer-1}{}}
\newlabel{lemma-compare-cosk0@cref}{{[originallemma][0][]01GC}{[1][0][]source}{}{}{}}
\makeatother
\crefname{originallemma}{original lemma}{original lemmas}
\Crefname{originallemma}{Original lemma}{Original lemmas}
\makeatletter
\newlabel{simplicial-lemma-section}{{01A6}{source}{}{sourcepointer-2}{}}
\newlabel{simplicial-lemma-section@cref}{{[originallemma][0][]01A6}{[1][0][]source}{}{}{}}
\makeatother
\crefname{originallemma}{original lemma}{original lemmas}
\Crefname{originallemma}{Original lemma}{Original lemmas}
\makeatletter
\newlabel{simplicial-lemma-trivial-kan-homotopy}{{08NS}{source}{}{sourcepointer-3}{}}
\newlabel{simplicial-lemma-trivial-kan-homotopy@cref}{{[originallemma][0][]08NS}{[1][0][]source}{}{}{}}
\makeatother
\title{Homology isomorphism from a level-n sheaf-surjective coskeleton map}
\author{The Stacks Project Authors}
\date{stacks-project:hypercovering.tex:local-snapshot-20261012; GNU FDL 1.2 or later}
\begin{document}
\maketitle
\noindent\textit{Editorial note.} Statement span is lines 879-893 of sources/stacks-hypercovering/pinned-original-source.tex; proof span is lines 895-932, the complete author proof (36 body lines).\par\medskip
\noindent\textit{Editorial note.} Navigation stubs: 3 cited-but-not-included labels are supplied as non-load-bearing pointers with their pinned tags where the snapshot tag index has them.\par\medskip
\noindent\textit{Editorial note.} extra\_packages is empty on purpose: the retained preamble already loads xy with options and a second declaration clashes in the assembler wrapper.\par\medskip
\noindent\textit{Editorial note.} Text is retained verbatim, including the author's own typos and placeholders.\par\medskip
\noindent\textit{Editorial note.} Queue index 13 of queue\_hypercovering.json (pb=1406, stmt 879-893, proof 895-932).\par\medskip
\noindent\textit{Editorial note.} stacks-hypercovering wave-1 (hc01): fresh adjudication accept (2026-10-14); pool per standing fallback order (hypercovering -> pic -> spaces-more-cohomology) after md01 closed the moduli pool; hypercovering.tex 53 labels/34 with proof judged 26 accept / 8 reject / 14 pb=0 env-dup rows excluded; spans cite-free and reference/slogan/history-free so no patchenv needed.\par\medskip
\noindent\textit{Allowed standard prerequisites (editorial).}\begin{itemize}
\item Explicit prerequisite boundary. The cited inputs are prerequisites; the argument in the retained span is not replaced by any cited result.
\end{itemize}\medskip
\ifdefined\hypertarget\hypertarget{sourcepointer-1}{}\fi
\noindent\textit{Source reference.} Original lemma 01GC: sources/stacks-hypercovering/pinned-original-source.tex, source lines 808--824. Cited by the retained statement or proof but not part of the collected spans; supplied only as a navigation pointer. The argument does not depend on its content.\par\medskip
\ifdefined\hypertarget\hypertarget{sourcepointer-2}{}\fi
\noindent\textit{Source reference.} Original lemma 01A6: sources/stacks-hypercovering/pinned-original-source.tex, source lines 932--932. Cited by the retained statement or proof but not part of the collected spans; supplied only as a navigation pointer. The argument does not depend on its content.\par\medskip
\ifdefined\hypertarget\hypertarget{sourcepointer-3}{}\fi
\noindent\textit{Source reference.} Original lemma 08NS: sources/stacks-hypercovering/pinned-original-source.tex, source lines 932--932. Cited by the retained statement or proof but not part of the collected spans; supplied only as a navigation pointer. The argument does not depend on its content.\par\medskip
\begin{lemma}
\label{lemma-acyclicity}
Let $\mathcal{C}$ be a site.
Let $f : L \to K$ be a morphism of
simplicial objects of $\textit{PSh}(\mathcal{C})$.
Let $n \geq 0$ be an integer.
Assume that
\begin{enumerate}
\item For $i < n$ the morphism $L_i \to K_i$ is an isomorphism.
\item The morphism $L_n \to K_n$ is surjective after sheafification.
\item The canonical map $L \to \text{cosk}_n \text{sk}_n L$ is an isomorphism.
\item The canonical map $K \to \text{cosk}_n \text{sk}_n K$ is an isomorphism.
\end{enumerate}
Then $H_i(f) : H_i(L) \to H_i(K)$ is an isomorphism.
\end{lemma}

\begin{proof}
This proof is exactly the same as the proof of
Lemma \ref{lemma-compare-cosk0} above. Namely,
we first let $K_n' \subset K_n$ be the sub presheaf
which is the image of the map $L_n \to K_n$. Assumption
(2) means that the sheafification of $K_n'$ is equal to
the sheafification of $K_n$. Moreover, since $L_i = K_i$
for all $i < n$ we see that get an $n$-truncated
simplicial presheaf $U$ by taking
$U_0 = L_0 = K_0, \ldots, U_{n - 1} = L_{n - 1} = K_{n - 1}, U_n = K'_n$.
Denote $K' = \text{cosk}_n U$, a simplicial presheaf.
Because we can construct $K'_m$ as a finite limit, and
since sheafification is exact, we see that
$(K'_m)^\# = K_m$. In other words, $(K')^\# = K^\#$.
We conclude, by exactness of sheafification once more,
that $H_i(K) = H_i(K')$. Thus it suffices to prove the lemma
for the morphism $L \to K'$, in other words, we may
assume that $L_n \to K_n$ is a surjective morphism
of {\it presheaves}!

\medskip\noindent
In this case, for any object $U$ of $\mathcal{C}$ we
see that the morphism of simplicial sets
$$
L(U) \longrightarrow K(U)
$$
satisfies all the assumptions of
Simplicial, Lemma \ref{simplicial-lemma-section}.
Hence it is a trivial Kan fibration. In particular it is
a homotopy equivalence
(Simplicial, Lemma \ref{simplicial-lemma-trivial-kan-homotopy}).
Thus
$$
\mathbf{Z}_L(U) \longrightarrow \mathbf{Z}_K(U)
$$
is a homotopy equivalence too. This for all $U$.
The result follows.
\end{proof}

\end{document}
